\section{The spectral geometry of hyperbolic and spherical spaces}
\begin{frame}{The Laplace-Beltrami operator}
	\begin{enumerate}
		\begin{definition}[Laplace-Beltrami operator]
			a generalization of the Laplace operator to more general spaces
			\begin{align*}
				\Delta f \coloneq \divergence \nabla f.
			\end{align*}
		\end{definition}
		\item The spectrum of this operator is a subset of $[0, \infty)$. We call it the \textcolor{blue}{\emph{spectrum}} of the manifold.
		      \pause
		\item The spectrum determines many things about the space (like its volume).
		\item When two distinct (up to isomorphism) manifolds share a spectrum, they for an \textcolor{red}{\emph{isospectral pair}}.
	\end{enumerate}
\end{frame}

\begin{comment}
\begin{frame}{Can isospectral pairs approximate your mom}
	\begin{itemize}
		\item One of the problems the paper studies is finding volume-maximizing isospectral pairs of spherical forms (i.e. the spaces we are interested in).
		\item Isospectral pairs share the order of their homotopy groups.
		\item The paper classifies spherical forms with homotopy group of order $<24$.
		\item It follows that any such isospectral pairs (with group order $<24$) are lens spaces.
	\end{itemize}
\end{frame}

\begin{frame}{Combining lenses into larger spaces}
	\begin{itemize}
		\item We can combine lens spaces by concatenating their lists of indices.
		\item Together with a computational search, the paper builds a table of the solutions for a bunch of dimensions. The general case is still an open question.
	\end{itemize}
\end{frame}
\end{comment}

{
\usebackgroundtemplate{
	\begin{tikzpicture}
		\clip (0,0) rectangle (\paperwidth,\paperheight);
		\fill[color=silver] (\paperwidth-10pt,0) rectangle (\paperwidth,\paperheight);
		% Added
		\fill[color=silver](0,0) rectangle (10pt,\paperheight);
	\end{tikzpicture}
}

\begin{frame}{The silver lining}
	\begin{itemize}
		\item Isospectral pairs are intriguing in their own right, but they cannot occur in dimension $3$.
		\item We can thus attempt to infer the shape of our universe based on its spectrum.
	\end{itemize}
\end{frame}
}
